\section{Difficulties and resolutions}
\label{sec:difficulties}

Every environment of this kind embodies dozens of small decisions made
under failure. Hiding them would make the result irreproducible: a
reader rebuilding the environment would meet the same failures with no
record of the fixes. This section catalogues each difficulty in the
order we met it --- contact, geometry and control, then collection,
training and tooling --- with its symptom, diagnosis and resolution, and
closes with what remains unjustified.

\subsection{Contact grasping}
\label{sec:diff-grasp}

The first cube grasps failed roughly one attempt in three: the jaws
closed beside the cube, or clipped its edge and spun it away. The
diagnosis was not contact physics but tracking: the differential solver
carries an azimuth-dependent steady-state error of roughly
\SI{8}{\milli\metre}, so a script that positions the gripper perfectly
in its own frame places the real pinch centre off the cube. The earlier
payload, a ten-centimetre prism, had masked this error entirely.

The obvious fix --- a proportional servo centring the pinch during the
whole approach --- made matters worse, dropping success to zero. The
pipeline lags its target by several ticks, so the servo integrated stale
error into wind-up and swept the gripper past the cube. The working
form holds the approach for a convergence dwell, then corrects with a
low gain under a hard magnitude clip, and freezes the correction the
moment the jaws start closing.

The decisive improvement came from visual review of failure renders:
the fingertips hovered just above the cube's top face at the commanded
grasp depth. Commanding the pinch \emph{below} the cube's centre ---
lower than any geometrically sensible target --- lets the solver's
vertical undershoot land the pinch on the cube. With the dwell, the
clipped servo and the deepened grasp, single-arm grasping became
fully reliable.

\subsection{Place-phase dragging}
\label{sec:diff-drag}

Relay episodes then failed at the \emph{second} place: lengthening the
settle dwell --- normally the safe change --- made success \emph{worse}.
The hold servo was the culprit: after the cube touched down but before
the jaws opened, the servo kept correcting through the pipeline's lag
and dragged the placed cube off its target; more dwell meant more
dragging. The resolution freezes the hold correction once a previously
lifted cube is detected back at rest while still gripped. The general
lesson: any closed-loop correction acting through a laggy actuation
path needs an explicit end-of-life rule, not just a gain.

\subsection{The fixed simple-mode scenario}
\label{sec:diff-simpleseed}

Simple mode stakes an entire dataset on repeated demonstrations of one
fixed scenario (Section~\ref{sec:experiments}), so that scenario must be
one the oracle solves at a high per-attempt rate. The relay's default
scenario was not: the teleoperation oracle completed it only about once
in four attempts, every failure in the place phase, and collection
aborted at its minimum-yield floor before writing a usable dataset. The
single-arm task, whose default scenario the oracle solves almost always,
never showed the problem, which is what made the cause easy to
misattribute.

Two oracles turned the diagnosis into a controlled experiment. The
lag-free analytic oracle solved the very same scenario on every attempt,
while the real-time teleoperation oracle --- the one collection actually
uses --- solved it only intermittently. The deficit therefore lay in the
teleoperation pipeline's tracking lag on that particular contact
configuration, not in the contact script or the scene geometry, both of
which the two oracles share. A confirming negative result followed:
raising the place-phase servo gain, the change the script-tuning
hypothesis predicts, moved the rate not at all. The resolution is to
choose the fixed scenario by a short seed search rather than accept an
arbitrary default; the selected relay scenario is then solved on every
attempt. The lesson is that a fixed-scenario dataset inherits the
oracle's worst case, not its average, so the scenario must be measured
and selected, never assumed.

\subsection{Workspace envelope versus the twin table}
\label{sec:diff-envelope}

Teleoperation targets pass through an analytic workspace envelope whose
floor protects the real rig's table surface. The twin's table top sits
about \SI{34}{\milli\metre} below the plane that floor assumes, so
table-level grasp targets in simulation were silently projected upward
--- the gripper stopped short of the cube with no error anywhere. The
fix scopes an envelope-floor override to the simulated collection scene
only; the real stack's tuned floor is untouched. The implication worth
recording for the real rig: whenever the physical working surface
differs from the height the envelope was tuned at, grasps near the
surface will fail in exactly this silent way.

\subsection{The abandoned mid-air hand-over}
\label{sec:diff-handover}

The original bimanual task held the object in the air while the second
arm grasped it. Two failure modes resisted all tuning. First, a
parallel-jaw grasp leaves the object's yaw about the vertical axis
unconstrained; during the carry the object rotated freely, so the
second gripper's approach met an unpredictable cross-section. Second,
two gripper bodies converging on one small object collided with each
other. Both problems vanish when the object is staged on the table
between grasps, which is why the relay of Section~\ref{sec:tasks}
exists. The lesson generalises: with parallel jaws and no wrist
redundancy, prefer task decompositions in which every grasp starts from
a fully constrained object pose.

\subsection{Renderer cache}
\label{sec:diff-renderer}

The twin cached one offscreen renderer keyed by image width alone.
Requesting two camera resolutions silently returned buffers of the
wrong height for one of them. The cache is now keyed by the full
resolution. A one-line bug, but it cost a debugging session because the
symptom (distorted frames in one stream) appeared far from the cause.

\subsection{Non-monotonic training and checkpoint selection}
\label{sec:diff-checkpoint}

A from-scratch imitation policy does not improve monotonically with
training step. On the single-arm task one intermediate checkpoint placed
the cube on every trial, yet several of its neighbours --- both earlier
and later --- scored zero, closing the jaws on empty air with the cube
never moving. Reading only the final checkpoint would have reported the
whole method as a failure. This is a property of the optimisation rather
than a bug to be removed, so the resolution treats checkpoint choice as
part of the method: every saved checkpoint rolls out on a held-out
validation seed set and the best advances to evaluation
(Section~\ref{sec:experiments}), ties broken towards the later step. The
price is one validation sweep per run; the alternative is mistaking a
training-step artefact for a verdict on the policy.

\subsection{Training-stack pitfalls}
\label{sec:diff-tooling}

For reproducibility, the pitfalls met wiring three policy families to
locally collected datasets, none of which produced a helpful error
message at the point of failure:

\begin{itemize}
  \item \textbf{Normalisation lives outside the policy.} The framework's
    policies do not normalise inside their action-selection call;
    normalisation is a separate processor pipeline saved alongside the
    checkpoint. An evaluation harness that loads only the policy weights
    silently evaluates an unnormalised model. Ours loads the
    checkpoint's own processor pipelines and passes every observation
    and action through them.
  \item \textbf{Observation construction must match the dataloader.}
    Evaluation images must be channel-first floating point in the unit
    range, keyed exactly as the dataset features were; we reuse the
    framework's own observation-conversion helper rather than
    reimplementing it.
  \item \textbf{Output directories must not pre-exist.} The trainer
    refuses to write into an existing run directory; resumable scripts
    must reuse finished cells and clear partial ones.
  \item \textbf{Video decoding.} Dataset videos are encoded with a
    modern codec; the default decoding backend loads system libraries
    that a sudo-less machine may lack. The pure-Python backend bundles
    its own decoder and is passed explicitly everywhere.
  \item \textbf{Gated model access.} The vision-language policy's
    tokenizer lives in a licence-gated repository; an unauthenticated
    machine fails at download time deep inside training. Our scripts
    check access in preflight and either skip the cell (smoke test) or
    abort before any compute is spent (long run).
  \item \textbf{Model sizing.} Full finetuning of the
    four-billion-parameter vision-language policy does not fit a
    \SI{24}{\gibi\byte} GPU once optimiser state is counted; low-rank
    adaptation of the published base fits comfortably and is the
    configuration our long-run script uses.
  \item \textbf{Multi-camera activation memory.} The from-scratch
    diffusion policy builds a separate convolutional encoder per camera
    and, by default, encodes frames at collection resolution; three
    full-resolution streams exhaust the same \SI{24}{\gibi\byte} GPU at
    the batch size the other policies train at. This is an activation, not
    a parameter, cost, so low-rank adaptation cannot help; downsampling
    the camera observations before the encoders restores ample headroom
    and costs nothing on a task whose control signal is already
    low-dimensional.
  \item \textbf{Fixed observation keys in a pretrained policy.} The
    vision-language base names its camera inputs by a fixed convention; a
    locally collected dataset whose streams are named differently aborts
    training with a feature mismatch before the first step. A rename map,
    baked into the saved observation processor so that evaluation inherits
    it unchanged, aligns the dataset to the policy without altering
    either.
\end{itemize}

\subsection{What remains unjustified}
\label{sec:diff-unjustified}

The contact constants (friction, impedance ratio, pad geometry), the
payload mass, the servo gains, clips and dwell durations, and the
oracle-gate thresholds beyond their measured bands are all tuned rather
than derived; each is flagged where introduced. They are stable across
the seeds exercised so far, and the oracle gate exists precisely to
catch a future regression in any of them before it can poison a
training run.

\subsection{A control rate that must divide the integrator}
\label{sec:diff-rate}
The control rate is a parameter of the simulated environment rather than
a constant, and a requested rate is refused unless it divides the
physics rate exactly. The refusal matters more than the flexibility. A
control period that is not a whole number of integrator steps leaves
recorded timestamps drifting against the motion they label, a fraction
of a step per tick, accumulating silently across an episode; the
recording still looks well formed, and nothing in the pipeline
announces the discrepancy. Refusing the rate at the point it is asked
for turns a slow corruption of every downstream timestamp into an error
message.
